\documentclass[../../main.tex]{subfiles}
\begin{document}

{\bf [解]}
題意の関係式
\begin{align}
 2^{n-1}\alpha=a_n+b_n \label{eq:1}
\end{align}
の両辺に$2$をかけて，
\begin{align}
2^n\alpha=2a_n+2b_n=a_{n+1}+b_{n+1} \quad \cdots \text{①} \label{eq:2}
\end{align}
が$a_n, b_n$についての漸化式である．
題意より$b_n$の大きさが$n$の偶奇で異なることから場合分けして考える．
以下 $k\in\mathbb{N}$ とする．

\subparagraph{$1^\circ$ $n=2k$の時}
この時は題意より$1<2b_n<2$だから，\cref{eq:2}は
\begin{align}
\begin{cases}\displaystyle
a_{2k+1}=2a_{2k}+1 \\
b_{2k+1}=2b_{2k}-1
\end{cases}
\quad \cdots \text{②} \label{eq:3}
\end{align}
となる．

\subparagraph{$2^\circ$ $n=2k-1$の時}
この時は題意より$0\le2b_n<1$だから，\cref{eq:2}は
\begin{align}
\begin{cases}\displaystyle
a_{2k}=2a_{2k-1} \\
b_{2k}=2b_{2k-1}
\end{cases}
\quad \cdots \text{③} \label{eq:4}
\end{align}
となる．

以上\cref{eq:3,eq:4}からまず$n$が奇数の時の$b_n$について
\begin{align}
 b_{2k+1}=4b_{2k-1}-1 \quad\therefore\ b_{2k+1}-\dfrac{1}{3}=4\left(b_{2k-1}-\dfrac{1}{3}\right) \label{eq:5}
\end{align}
となる．$b_1=\alpha$ だから\cref{eq:5}をくり返し用いて，
\begin{align}
b_{2k-1}=4^{k-1}\left(\alpha-\dfrac{1}{3}\right)+\dfrac{1}{3}  \label{eq:6}
\end{align}
を得る．\cref{eq:3}より$n$が偶数の時の$b_n$について
\begin{align}
b_{2k}=2\cdot4^{k-1}\left(\alpha-\dfrac{1}{3}\right)+\dfrac{2}{3} \label{eq:7}
\end{align}
となる．

さて，題意より $0\le b_n<1$ が任意の$n$で成立するので，$\alpha=\dfrac{1}{3}$ であることが必要で，逆にこの時\cref{eq:6,eq:7}は
\begin{align}
b_{2k-1}=\dfrac{1}{3},\quad b_{2k}=\dfrac{2}{3} \label{eq:8}
\end{align}
となって条件を満たす．よって $\alpha=\dfrac{1}{3}$ である．

$a_n$は\cref{eq:8}を\cref{eq:1}に代入して
\begin{align}
\begin{split}
a_{2k-1}&=\dfrac{1}{3}\left(4^{k-1}-1\right) \\
a_{2k}&=\dfrac{2}{3}\left(4^{k-1}-1\right)
\end{split}
\end{align}
と求まる．

\newpage

\end{document}
